% supplement/supplement.tex: the Supplementary Information.
% Build it from the Final Paper folder, not from supplement/:
%   python3 tools/make_numbers.py, then latexmk -pdf supplement/supplement.tex
% All paths below are relative to the Final Paper folder. Numbers come only from
% the macros in generated/numbers.tex. The tables and figures that the pipeline
% writes are in supplement/tables/ and supplement/figures/; input them with
% \input{supplement/tables/<name>} and \includegraphics{<name>}.
\documentclass[11pt,a4paper]{article}
% Allow the PDF 1.7 figures that the pipeline writes to be included without a warning.
\pdfminorversion=7

\usepackage[T1]{fontenc}
\usepackage[utf8]{inputenc}
\usepackage{lmodern}
\usepackage[margin=2.5cm]{geometry}
\usepackage{amsmath,amssymb}
\usepackage{graphicx}
\usepackage{booktabs}
\usepackage{array}
\usepackage{multirow}
\usepackage{longtable}
\usepackage{float}
\usepackage[labelfont=bf]{caption}
\usepackage[numbers,super,sort&compress]{natbib}
\usepackage[hidelinks]{hyperref}
% Long URLs in the reference list may break at any character.
\usepackage{xurl}
% The middle dot in the name DALL\textperiodcentered E: Latin Modern gives it wide side bearings, so
% close them up to print the name as one word. The paper uses the middle dot nowhere else. This
% waits for the start of the document because hyperref redefines the command when it loads.
\AtBeginDocument{%
  \let\origtextperiodcentered\textperiodcentered
  \renewcommand{\textperiodcentered}{\kern-0.08em\origtextperiodcentered\kern-0.08em}}

\graphicspath{{supplement/figures/}{figures/}}
% Tables and figures are numbered 1, 2, ... and their labels are named in full:
% "Supplementary Table 1", "Supplementary Figure 1". Running text, here and in the main
% text, cites them as "Supplementary Table 1" and "Supplementary Fig. 1".
\renewcommand{\tablename}{Supplementary Table}
\renewcommand{\figurename}{Supplementary Figure}
% Supplementary Notes are numbered 1, 2, ... in the order the main text's Results cite them.
\newcounter{suppnote}
\newcommand{\suppnote}[1]{\refstepcounter{suppnote}\subsection*{Supplementary Note \thesuppnote: #1}\addcontentsline{toc}{subsection}{Supplementary Note \thesuppnote: #1}}

% The numbers the paper prints, one macro per number.
\input{generated/numbers}

\title{Supplementary Information for\\
  \textit{Beyond the Original Sin: Lay Evaluations of Creative Contribution in AI-Assisted Artworks}}
\author{David G. Kamper\textsuperscript{1,2,*}, Alice Lin\textsuperscript{1} and Keith J. Holyoak\textsuperscript{1}\\[1ex]
  \small\textsuperscript{1}Department of Psychology, University of California, Los Angeles, Los Angeles, CA, USA\\
  \small\textsuperscript{2}University of Chicago Law School, Chicago, IL, USA\\
  \small\textsuperscript{*}Corresponding author: davidgkamper@ucla.edu}
\date{}

\begin{document}

\maketitle

\noindent This document has four parts. Section~\ref{sec:supp-verdicts} gives Supplementary Tables~1 and~2: every registered test with its status and verdict, and the AI label's effect on second works. Section~\ref{sec:supp-notes} holds Supplementary Notes~1 to~7, in the order in which the main text's Results cite them. Section~\ref{sec:supp-methods}, the Supplementary Methods, gives the detail behind the main text's Methods. Sections~\ref{sec:suppA-registrations} to~\ref{sec:suppB-crossstudy} give the registrations, then each study's materials, procedure, sample and remaining results, then the text measures shared by all five studies and the exploratory cross-study analyses. Each study's section begins by naming the notes and tables that hold its other results.

\tableofcontents
\clearpage

% Supplementary Tables 1 and 2, written by tools/make_tables.py (run it after
% tools/make_numbers.py). They come first so that they carry those numbers.
\section{Registered tests and their verdicts}\label{sec:supp-verdicts}
\input{supplement/sections/verdicts}
\input{supplement/tables/supp_table1_verdicts}
\input{supplement/tables/supp_table2_label_effects}

% The main text's Results point to these notes by number; keep this order.
\section{Supplementary Notes}\label{sec:supp-notes}
\input{supplement/sections/note1_study1_h3_h4}
\input{supplement/sections/note2_study1_robustness_belief}
\input{supplement/sections/note3_study3_detail}
\input{supplement/sections/note4_image_check}
\input{supplement/sections/note5_study5_editing_time}
\input{supplement/sections/note6_pooled_second_works}
\input{supplement/sections/note7_explanations}

\input{supplement/sections/supp_methods}

\section{Registrations}
\input{supplement/sections/registrations}

\section{Study 1}
\input{supplement/sections/study1}

\section{Study 2}
\input{supplement/sections/study2}

\section{Study 3}
\input{supplement/sections/study3}

\section{Study 4}
\input{supplement/sections/study4}

\section{Study 5}
\input{supplement/sections/study5}

\section{Text measures shared by all five studies}
\input{supplement/sections/text_measures}

\section{Exploratory cross-study analyses}
\input{supplement/sections/crossstudy}

\bibliographystyle{naturemag}
\bibliography{references}

\end{document}
